% !TeX root = ../main.tex
\section{Coupled elastic verification}
\label{sec:mandel-verification}

\subsection{Boundary-value problem}
\label{sec:mandel-boundary-value-problem}

The Mandel problem considers a saturated plane-strain specimen compressed
between impermeable, frictionless rigid platens.  Water drains through the
lateral faces \citep{chengdetournay1988}.  Let \(X\) and \(Y\) denote the
horizontal and vertical reference coordinates, and let \(L\) and \(H\) be the
quarter-domain width and height, respectively.  Symmetry reduces the
calculation to
\begin{equation}
 \Omega_0=\left\{(X,Y):0\leq X\leq L,\ 0\leq Y\leq H\right\},
 \qquad
 L=1\ \mathrm{m},
 \qquad
 H=0.1\ \mathrm{m}.
 \label{eq:mandel-quarter-domain}
\end{equation}
The mechanical symmetry conditions and drained pressure condition are
\begin{subequations}
 \label{eq:mandel-boundary-conditions}
 \begin{align}
  u_x(0,Y,t)&=0,
  &u_y(X,0,t)&=0,
  \label{eq:mandel-symmetry-boundaries}\\
  p(L,Y,t)&=0,
  &\mbf W_f\mathbin{\cdot}\mbf N&=0
  \quad\text{on the remaining fluid boundaries}.
  \label{eq:mandel-fluid-boundaries}
 \end{align}
\end{subequations}
Here \(u_x\) and \(u_y\) are the components of \(\mbf u\), and \(\mbf N\) is
the outward unit normal to the reference boundary.
The top displacement is spatially uniform and follows the analytical
constant-load solution of \citet{chengdetournay1988}.  The computed vertical
nominal-stress resultant then provides an independent check of the applied
\(100\ \mathrm{kPa}\) compression.  The right face is mechanically
traction-free, and the prescribed-normal-displacement boundaries are
frictionless.  The initial material state is
\(\mbf u=\mbf0\), \(p=0\), and \(\rho_s=\rho_{s0}\); the local constitutive
state then has \(\bar\rho_s=\bar\rho_{s0}\), \(\phi_s=\phi_{s0}\), and
\(\bar J=1\).  The first time step applies the analytical platen
displacement, after which its time history follows
\eqref{eq:mandel-vertical-displacement-series}.  Because the vertical
displacement is imposed and is linear in \(Y\), its profile is a
boundary-condition reproduction check rather than an independent benchmark
observable.

The solid-mass residual in \eqref{eq:solid-mass-weak-residual} evolves
\(\rho_s\) throughout the calculation.  At each integration point, the
mineral equation determines intrinsic density, the density ratio gives
\(\phi_s\), and \eqref{eq:poroplastic-biot-correction} gives \(B\) with
\(a^p=1\).  Water density follows
\eqref{eq:barotropic-water-density}.  The small load keeps the coefficient
near its reference value while testing the nonlinear implementation.

The material parameters are listed in \cref{tab:mandel-parameters}.
\begin{table}[ht]
 \centering
 \caption{Solid--water Mandel parameters.}
 \label{tab:mandel-parameters}
 \begin{tabular}{lr}
  \toprule
  Parameter & Value \\
  \midrule
  Drained skeleton bulk modulus \(K\) & \(1.0\ \mathrm{GPa}\) \\
  Skeleton shear modulus \(G\) & \(0.75\ \mathrm{GPa}\) \\
  Mineral bulk modulus \(K_s\) & \(2.5\ \mathrm{GPa}\) \\
  Water bulk modulus \(K_f\) & \(8.0\ \mathrm{GPa}\) \\
  Initial water volume fraction \(\phi_{f0}\) & \(0.1\) \\
  Permeability \(\kappa\) & \(1.5\times10^{-12}\ \mathrm{m^2}\) \\
  Water viscosity \(\mu_f\) & \(10^{-3}\ \mathrm{Pa\,s}\) \\
  Applied compressive stress \(P_0\) & \(100\ \mathrm{kPa}\) \\
  \bottomrule
 \end{tabular}
\end{table}

The reference coefficient from
\eqref{eq:poroplastic-biot-correction} is \(B_0=0.6\).


\subsection{Pressure and displacement}
\label{sec:mandel-numerical-comparison}

The linearization and analytical series are given in
Appendix~\ref{sec:mandel-analytical-solution}.  The storage coefficient is
obtained from the matched logarithmic mineral law and the water law.  It
recovers the classical single-mineral storage relation, so the benchmark
tests the small-strain limit of the finite-deformation implementation.

The error and refinement study uses a \(40\)-by-\(4\) QUAD9 mesh with
uniform backward-Euler steps \(\Delta t=0.002\ \mathrm{s}\).  The spatial
profiles use the same mesh with an early-refined sequence beginning at
\(0.001\ \mathrm{s}\) and increasing to \(0.016\ \mathrm{s}\).
\Cref{fig:mandel-pressure-profiles} compares pressure along
\(Y=0.05\ \mathrm{m}\) at several times.  The solution reproduces the
pressure rise near the undrained center, followed by decay as fluid leaves
through the lateral boundary.

\begin{figure}[ht]
 \centering
 \figureinput{mandel_pressure_profiles.pgf}
 \caption{Spatial water-pressure profiles during the small-load Mandel test.
 Solid curves are the analytical series and open circles are the Q2/Q1
 finite-element solution.  The dashed line marks the instantaneous undrained
 pressure.  The early profile exceeds this level near the center before
 drainage reduces the pressure throughout the specimen.}
 \label{fig:mandel-pressure-profiles}
\end{figure}

\Cref{fig:mandel-displacements} compares the displacement profiles at the
same times.  Horizontal displacement is sampled along the specimen mid-height and
provides an independent check of the coupled mechanical response.  Vertical
displacement is sampled on the drained side; its agreement confirms
reproduction of the imposed analytical platen history.

\begin{figure}[ht]
 \centering
 \figureinput{mandel_displacements.pgf}
 \caption{Horizontal and vertical displacement profiles in the elastic Mandel
 verification.  Solid curves are the analytical series and open circles are
 finite-element values.  Horizontal displacement is an independent benchmark
 observable; vertical displacement checks the prescribed platen motion.}
 \label{fig:mandel-displacements}
\end{figure}

The verification evaluates pressure over the complete history at three
horizontal positions and checks the recovered nominal load.  Separate mesh
and time-step comparisons measure numerical sensitivity.  Successive temporal
refinements reduce the pressure difference, while the mesh comparison shows
smaller spatial sensitivity on the tested grids.  These pairwise comparisons
measure sensitivity over the reported transient; they do not establish an
asymptotic convergence order.

The assembled-Jacobian checks described in
Sec.~\ref{sec:weak-residuals-and-spaces} also pass for this problem.
The mineral EOS residual and referential solid-mass drift are recorded
separately.  The accompanying validation records give numerical errors,
tolerances, and sampling definitions.
\FloatBarrier

\subsection{Finite-deformation continuation}
\label{sec:mandel-finite-deformation-continuation}

A second calculation uses the same domain, phase balances, and elastic
constitutive model at a larger imposed deformation.  The platen displacement
increases linearly to \(-0.1H\) over \(0.2\ \mathrm{s}\) and is then held.
The pressure, lateral displacement, nominal stress, and coefficient are
predictions of this displacement-controlled calculation.

\begin{figure}[ht]
 \centering
 \figureinput{mandel_finite_deformation.pgf}
 \caption{Finite-deformation continuation of the coupled elastic system.
 Left: imposed platen compression and predicted lateral expansion.
 Right: spatial average and range of the Biot coefficient.  The vertical
 dotted line marks the end of the displacement ramp.}
 \label{fig:mandel-finite-deformation}
\end{figure}

The coefficient varies across the specimen during compression.  It can
briefly exceed the reference value in the undrained interior, while the
largest reduction occurs near the drained boundary
(\cref{fig:mandel-normalized-biot-contours}).  Drainage during the hold
reduces lateral expansion and brings the coefficient toward a nearly
uniform value below its reference value.

\begin{figure}[ht]
 \centering
 \figureinput{mandel_normalized_biot_contours.pgf}
 \caption{Normalized Biot coefficient during finite compression and drainage.
 The panels share a color scale and use solid reference coordinates.  The
 largest transient reduction occurs near the drained boundary.  The vertical
 plotting scale is enlarged to show the through-height field.}
 \label{fig:mandel-normalized-biot-contours}
\end{figure}

The coefficient identity and mineral EOS remain satisfied during this
continuation.  The solid-mass diagnostic measures the discretization error
of the weak chain-rule balance, and its drift decreases under time
refinement.  Mesh and time-step comparisons also check pressure, lateral
displacement, and the spatially averaged coefficient.  This calculation tests the elastic model at finite deformation.  Plastic
history is examined separately in the material-point calculations of
Sec.~\ref{sec:poroplastic-demonstration}.
\FloatBarrier

\subsection{Scope and reproducibility}
\label{sec:verification-scope}

The source archive contains the MOOSE materials and kernels, input decks,
independent constitutive checks, analytical reference calculations, curated
results, and figure-generation scripts.  The equation-to-object and theory
maps link each test to its governing equations and assumptions.  The
calculations verify the implementation for the stated constitutive laws;
comparison with material data remains a separate task.
